\section{Trabalho relacionado}
\label{sec:related}

\subsection{O dataset TON\_IoT e o paper de Alsaedi 2020}

O \textbf{TON\_IoT}~\cite{alsaedi2020toniot,moustafa2021toniot} foi
construído pela equipe da UNSW Canberra Cyber e cobre três camadas
distintas do tráfego de uma rede IoT: tráfego de rede, eventos de
sistema operacional e \textbf{telemetria de sensores}. Este relatório
foca exclusivamente na camada de telemetria, que contém sete bases:
\textbf{Fridge}, \textbf{Garage\_Door}, \textbf{GPS\_Tracker},
\textbf{Modbus}, \textbf{Motion\_Light}, \textbf{Thermostat} e
\textbf{Weather}, totalizando cerca de 260 mil linhas, mais uma base
\textbf{combined} (concatenação das sete, usada nas Tabelas 12--13 do
paper). Cada amostra é rotulada como \texttt{normal} ou como um dos
sete tipos de ataque (\texttt{backdoor}, \texttt{ddos},
\texttt{injection}, \texttt{password}, \texttt{ransomware},
\texttt{scanning}, \texttt{xss}; nem toda base contém todos os tipos).

O paper de Alsaedi et al.\ avalia oito \emph{baselines} clássicos
nessa telemetria sob três tarefas: (i) classificação binária
\emph{per-device} (Tabelas 10--11 do paper), (ii) classificação binária
sobre o \emph{combined} (Tabela 12), e (iii) classificação
multi-classe sobre o \emph{combined} (Tabela 13). O protocolo,
descrito na Seção VI--A, é \emph{split} 80\%/20\% com validação
cruzada $k=4$ no \emph{train}, e a métrica reportada é F-score, embora
a chamada exata da função (\texttt{binary}, \texttt{macro}, ou
\texttt{weighted}) não seja explicitada. Tratamos essa ambiguidade na
Seção~\ref{sec:discussao}.

\subsection{Vazamento de dados e qualidade do TON\_IoT}
\label{sec:related_leak}

A comunidade já documentou vazamento de dados (\emph{data leakage} ---
quando uma \emph{feature} carrega informação do rótulo que não estaria
disponível em produção, inflando artificialmente a acurácia) no
TON\_IoT, porém sempre no \textbf{dataset de rede} e por
\emph{features} identificadoras. Shaikhanova et al.~\cite{shaikhanova2025audit}
mostram que o resultado de $99{,}97\%$ do estudo original incluía
endereços IP e números de porta como \emph{features}, atributos que
``introduzem vazamento e superestimam o desempenho real'', e propõem
descartar oito colunas textuais de alta cardinalidade. Dharini et
al.~\cite{dharini2026efficient} fazem a mesma crítica: os ataques foram
lançados de faixas de IP e portas pré-definidas, fazendo esses campos
agirem como identificadores de ataque em vez de indicadores de
comportamento. Raskovalov et al.~\cite{raskovalov2022rectification} vão
além e publicam uma versão corrigida (ToN-IoT-R), filtrando fluxos
rotulados como ataque que continham tráfego normal com o roteador e o
\emph{resolver} DNS. Todas essas correções operam sobre colunas
(\texttt{ssl\_subject}, \texttt{http\_uri}, IP, porta) que
\emph{inexistem} na telemetria de sensores --- portanto não alcançam o
vazamento que identificamos.

Na camada de \textbf{telemetria}, ao contrário, a acurácia
quase-perfeita é tratada como virtude, e não como sintoma. Sharrab et
al.~\cite{sharrab2025analysis} dedicam um artigo ao Garage\_Door e
reportam $100\%$ de acurácia para múltiplos classificadores, atribuindo
o resultado ao ``poder preditivo das \emph{features} mais
correlacionadas'' sem levantar qualquer alarme; Alotaibi e
Ilyas~\cite{alotaibi2023stacking} combinam seis dispositivos de
telemetria e reportam $98{,}64\%$ sem uma seção de limitações sobre a
acurácia. Benaddi et al.~\cite{benaddi2025near} reportam macro-F1 acima
de $0{,}986$ no TON\_IoT com um IDS comprimido (poda de \emph{features}
guiada por SHAP e distilação de conhecimento), tratando a acurácia
quase-perfeita como resultado a otimizar e não como sintoma a
investigar, sem qualquer análise de vazamento. A crítica de
heterogeneidade de Booij et al.~\cite{booij2022heterogeneity}, a mais
citada sobre o dataset, trata de generalização entre versões, não de
vazamento.

Nenhum desses trabalhos inspeciona os CSVs de telemetria no nível da
\emph{feature} categórica, e nenhum conecta a determinismo da
associação à sub-amostragem da classe \texttt{normal} nos arquivos
públicos. É exatamente esse o foco da nossa auditoria
(Seção~\ref{sec:dataset}): diagnosticamos um vazamento categórico
determinístico (Cramér's V $=1{,}000$) em \texttt{Fridge} e
\texttt{Garage\_Door} e o rastreamos até a sub-amostragem
$35\,\text{k} \rightarrow 15\,\text{k}$ da classe \texttt{normal},
um mecanismo distinto do vazamento de identificadores da camada de
rede e, até onde alcançamos na literatura, ainda não reportado na
telemetria.

\subsection{Família WiSARD de redes neurais sem peso}

Redes neurais \textbf{sem peso} (WNNs) substituem produtos
peso-entrada típicos das redes convencionais por consultas
$\mathcal{O}(1)$ em \emph{RAMs} indexadas por padrões binários de
entrada. A versão seminal --- \textbf{WiSARD}~\cite{aleksander1984wisard}
--- é o \emph{discriminador} mais simples possível: um conjunto de
RAMs $r_1, \dots, r_N$, cada uma endereçada por uma tupla de
$n$ \emph{bits} (\emph{tuple address}) sorteada (mapeamento
pseudo-aleatório) sobre a representação binária da entrada. Para cada
classe, treinar significa marcar com 1 as posições endereçadas pelas
amostras dessa classe; classificar significa contar quantas RAMs
respondem 1 para a entrada candidata e escolher a classe com maior
contagem.

A simplicidade tem custo: a representação binária da entrada precisa
preservar a estrutura métrica das \emph{features}. A solução clássica
é o \textbf{termômetro}, que codifica um valor real $x \in [a, b]$ em
$T$ \emph{bits}, com os primeiros $\lceil T \cdot (x-a)/(b-a) \rceil$
ligados. Variantes ajustadas aos dados (distributivo, gaussiano,
exponencial, logarítmico) adaptam a discretização à distribuição
empírica da \emph{feature}; avaliamos seis dessas variantes na
Seção~\ref{sec:metodologia}.

Os cinco modelos avaliados estendem o esqueleto WiSARD em direções
ortogonais:

\begin{description}[leftmargin=*,itemsep=0pt]
  \item[\textbf{WiSARD}~\cite{aleksander1984wisard}.] O modelo base.
        Hiperparâmetros principais: tamanho do endereço ($n$),
        termômetro e número de \emph{bits} por \emph{feature},
        \texttt{ignore\_zero} (descartar endereços nulos).
  \item[\textbf{ClusWiSARD}~\cite{cardoso2016cluswisard}.] Particiona
        cada classe em \emph{sub-discriminadores} via critério de
        similaridade (\texttt{min\_score}, \texttt{threshold},
        \texttt{limit}). Útil quando a classe tem sub-modos
        estatisticamente distintos.
  \item[\textbf{BloomWiSARD}~\cite{santiago2020bloomwisard}.]
        Substitui a RAM tradicional por um \emph{Bloom filter}, o que
        reduz drasticamente o consumo de memória ao custo de falsos
        positivos controlados.
  \item[\textbf{BTHOWeN}~\cite{susskind2022bthowen}.] Adota
        \emph{hashed lookup tables} de tamanho fixo
        (\texttt{filter\_entries}) com $k$ funções de hash
        (\texttt{num\_hashes}). Projetado para implementação em FPGA
        com latência sub-microsegundo.
  \item[\textbf{ULEEN}~\cite{susskind2023uleen}.] Adiciona uma fase
        adicional de \emph{learning rate} sobre as RAMs e um esquema
        de \emph{bleaching} treinável, alcançando F1 superior ao
        BTHOWeN em datasets de visão ao custo de mais iterações.
\end{description}

Os experimentos usam a biblioteca \texttt{wisardpkg}~\cite{wisardpkg}
estendida com o conjunto de termômetros ajustados descrito na
Seção~\ref{sec:metodologia}. Os modelos \emph{baseline} usam
\texttt{scikit-learn}~\cite{scikit-learn} para LR/LDA/kNN/RF/CART/NB/SVM
e \texttt{PyTorch} para o LSTM, sempre com semente fixa
(\texttt{random\_state=42}) em todos os nossos experimentos; o paper
não publica a semente que usou (ver \S{}\ref{sec:discussao}).
